% ============================================================
%  CAPÍTULO 9: CONCLUSIONES
% ============================================================
\chapter{Conclusiones}
\label{cap:conclusiones}

% ------------------------------------------------------------
\section{Síntesis de los resultados}
% ------------------------------------------------------------

Este Trabajo de Fin de Grado ha investigado los determinantes del gasto en seguros
sanitarios privados en España utilizando un panel de datos de las 17 comunidades
autónomas para el período 2016--2024. El marco teórico, basado en el modelo de
demanda de salud de Grossman y en la Ley de Little, predijo cuatro hipótesis
susceptibles de contraste empírico. Los resultados de las estimaciones econométricas
han permitido responderlas sistemáticamente.

La Tabla~\ref{tab:conclusiones_hip} resume el estado final de cada hipótesis a
la luz de la evidencia.

\begin{table}[htbp]
  \centering
  \caption{Estado final de las hipótesis de investigación}
  \label{tab:conclusiones_hip}
  \begin{threeparttable}
    \small
    \resizebox{\textwidth}{!}{%
    \begin{tabular}{clp{4cm}lc}
      \toprule
      \textbf{Hip.} & \textbf{Enunciado resumido}
        & \textbf{Resultado clave}
        & \textbf{Veredicto}
        & \textbf{Modelo} \\
      \midrule
      H1 & Bien de lujo ($\varepsilon_Y > 1$) &
        $\hat{\beta}_1 = 1{,}02^{***}$ (FE) &
        \textbf{Confirmada} & M1-FE \\
      H2 & Esperas $\uparrow \Rightarrow$ seguros $\uparrow$ (lineal) &
        $\hat{\gamma} = 0{,}0002$ (n.s.) &
        No confirmada & M1-FE \\
      H3 & Efecto diferido (retardos) &
        Wald $\chi^2(2) = 0{,}39$ ($p = 0{,}82$) &
        No confirmada & M2-FE \\
      H4 & Efecto umbral ($\tau = 100$ días) &
        $\hat{\delta} = 0{,}093^{**}$ ($p=0{,}017$) &
        \textbf{Confirmada} & M1-U \\
      H5 & Inercia del gasto &
        $\hat{\rho} = 0{,}210^{*}$ ($p=0{,}075$) &
        Parcialmente confirmada & M3-FE \\
      \bottomrule
    \end{tabular}}% fin resizebox
    \begin{tablenotes}[flushleft]
      \footnotesize
      \item $*$ $p < 0{,}10$;\quad $**$ $p < 0{,}05$;\quad $***$ $p < 0{,}01$.
      \item Test de Hausman ($\chi^2(4)=11{,}44$, $p=0{,}022$) indica M1-FE como
        modelo preferido.
    \end{tablenotes}
  \end{threeparttable}
\end{table}

Las cinco hipótesis han sido contrastadas con la evidencia empírica. Los resultados
más robustos son H1 (elasticidad renta significativa, $\hat{\beta} = 1{,}02$) y
H4 (efecto umbral significativo al 5\%). H5 (inercia del gasto) queda
parcialmente confirmada con significatividad al 10\% ($p = 0{,}075$). H2 (efecto
lineal de la espera) y H3 (efecto diferido con retardos) no quedan confirmadas:
el mecanismo de transmisión de las listas de espera hacia el gasto privado no
opera de forma lineal ni acumulativa, sino a través del efecto umbral.

% ------------------------------------------------------------
\section{Aportaciones principales del trabajo}
% ------------------------------------------------------------

Las contribuciones originales de este trabajo son las siguientes:

\begin{enumerate}
  \item \textbf{Primera estimación sistemática con datos de panel para las CCAA
    españolas} del período post-descentralización (2016--2024), incorporando el
    impacto de la pandemia COVID-19 como variable de control.
  \item \textbf{Identificación del efecto umbral}: la introducción de la variable
    dummy a 100 días demuestra que el efecto de las listas de espera opera de
    forma no lineal ---un incremento del 9{,}3\% en el gasto cuando se supera el umbral.
  \item \textbf{Documentación de la persistencia del gasto}: el modelo con gasto
    anterior (M3) revela inercia moderada en las decisiones de contratación,
    reflejando renovación de pólizas y formación de hábitos ($\hat{\rho}=0{,}21$,
    $p=0{,}075$).
  \item \textbf{Elasticidad renta como bien de lujo}: confirmación de que
    el seguro sanitario privado tiene elasticidad renta cercana o superior a uno
    ($1{,}02$ bajo FE, $1{,}54$ bajo RE), clasificándolo como bien de lujo a
    escala regional.
\end{enumerate}

% ------------------------------------------------------------
\section{Aportaciones para la política sanitaria}
% ------------------------------------------------------------

Los resultados del trabajo tienen implicaciones directas para el diseño de la
política sanitaria en España:

\begin{description}
  \item[Implicación I --- Gestión de listas de espera:] Reducir los tiempos de
    espera en consulta de especialistas por debajo del umbral de 100 días
    reduciría el gasto en seguros privados y reforzaría la cohesión del sistema
    sanitario público. El efecto umbral significativo (H4) implica que la
    prioridad debe focalizarse en las CCAA que superan este límite.

  \item[Implicación II --- Equidad:] La elasticidad renta superior a la unidad
    advierte de que el crecimiento económico, en ausencia de intervención política,
    tenderá a ampliar la segmentación entre usuarios del sistema público y del
    privado.

  \item[Implicación III --- Financiación autonómica:] Dado que las CCAA con mayor
    renta per cápita presentan mayor gasto en seguros privados, el sistema de
    financiación autonómica debería calibrar el esfuerzo sanitario considerando
    la demanda potencial sobre el sistema público, que es mayor en la medida en
    que la renta impulsa la salida de usuarios al sector privado.

  \item[Implicación IV --- Impacto del COVID-19:] El incremento del gasto en
    seguros privados durante 2020 (+17\% según M1-RE, +9\% según M1-FE) sugiere
    que las crisis sanitarias pueden acelerar la segmentación del sistema. Los
    sistemas sanitarios que dependen de la complementariedad público-privada
    deben planificar mecanismos de resiliencia que eviten que el colapso del
    sector público expulse a los hogares con capacidad de pago hacia
    alternativas privadas.
\end{description}

% ------------------------------------------------------------
\section{Reflexión final}
% ------------------------------------------------------------

La evidencia acumulada en este trabajo apunta a que el sistema sanitario español
funciona en la práctica como un \textbf{sistema de dos niveles}: un nivel público
universal al que acceden todos los ciudadanos, y un nivel privado complementario
que canaliza la demanda de los hogares con mayor renta y menor tolerancia a la
espera. Esta arquitectura dual no es ni inevitable ni necesariamente indeseable,
pero requiere una gestión activa para evitar que la calidad del primer nivel
se deteriore a medida que el segundo crece.

La sanidad como bien de lujo no es un destino irremediable, pero es una tendencia
que solo se puede reorientar mediante decisiones de política pública explícitas:
inversión en reducción de tiempos de espera, transparencia en la calidad de los
servicios públicos y financiación equitativa del SNS. La econometría no dicta esas
decisiones, pero sí proporciona las magnitudes que permiten evaluarlas con rigor.
